\documentclass[12pt]{article}

\usepackage{sbc-template}
\usepackage{graphicx,url}
\usepackage[utf8]{inputenc}
% citações autor-ano
\usepackage[authoryear]{natbib}

\sloppy

\title{Controle de um Robô com Linguagem Natural usando LBML e Aprendizado de Máquina}

\author{Guilherme Rosa\inst{1} }

\address{Programa de Pós-Graduação em Computação Aplicada -- UDESC\\
  Joinville, SC -- Brasil
  \email{guilherme.mr@udesc.br}
}

\begin{document}

\maketitle

\begin{abstract}
Este trabalho apresenta um sistema para controlar o robô LBot por meio de comandos em linguagem natural. A solução utiliza a linguagem de domínio específico LBML e um modelo de aprendizado de máquina para traduzir comandos em português.
\end{abstract}

\begin{resumo}
O artigo descreve um sistema que permite controlar o robô LBot usando comandos em linguagem natural. A abordagem emprega a linguagem LBML e um modelo de aprendizado de máquina para traduzir comandos em português, com avaliação de desempenho por métricas de tradução automática.
\end{resumo}

\section{Introdução}
Atualmente, dois campos que despertam grande interesse tanto na academia quanto na indústria são a Robótica e a Inteligência Artificial. O avanço nessas áreas tem proporcionado impactos significativos na sociedade, permitindo a criação de soluções inovadoras que já fazem parte do cotidiano.

O projeto apresentado neste artigo integra os campos de Robótica e Inteligência Artificial. Para isso, foi desenvolvido o robô LBot, cuja principal funcionalidade é a movimentação no espaço físico. A fim de possibilitar que esses movimentos sejam descritos de forma padronizada, foi criada uma linguagem de domínio específico (DSL), denominada LBot Movement Language (LBML).

Entretanto, um desafio recorrente em aplicações de robótica é a necessidade de que o usuário aprenda comandos técnicos ou linguagens específicas para interagir com o robô. Para contornar esse problema, este trabalho propõe o uso de um modelo de aprendizado de máquina capaz de traduzir comandos em português para a LBML, permitindo que usuários interajam com o LBot por meio de linguagem natural.

Diante desse contexto, este artigo apresenta a criação da LBML, o desenvolvimento do conjunto de dados de treinamento e a implementação de um modelo baseado em arquiteturas de transformers utilizando a biblioteca PyTorch, permitindo a tradução de comandos em português para a LBML. A combinação de linguagens de domínio e aprendizado de máquina demonstra como a interação entre humanos e robôs pode ser simplificada. Além disso, a utilização de comandos em linguagem natural torna o LBot mais acessível a usuários sem conhecimento técnico, potencializando aplicações práticas em educação, entretenimento e pesquisa em robótica.

\section{Revisão de Conceitos e Trabalhos Relacionados}

A interface entre a linguagem natural e as máquinas é um campo de pesquisa fundamental para a Robótica e a Inteligência Artificial (IA), sendo o núcleo deste trabalho o problema de Semantic Parsing (Análise Semântica). Esta seção explora os conceitos centrais e os trabalhos de ponta que pavimentaram o caminho para a tradução eficiente de comandos em linguagem natural para linguagens formais e de domínio específico.

\subsection{Semantic Parsing e Linguagens de Domínio Específico (DSL)}

O Semantic Parsing é a tarefa de traduzir frases em linguagem natural (LN) em representações formais e executáveis, como código de programação, expressões lógicas ou queries de banco de dados \citep{Sharma2025}. Embora o foco tradicional recaia sobre a conversão de perguntas para Structured Query Language (SQL) (o chamado NL-to-SQL), o conceito se estende a qualquer linguagem de máquina que defina ações em um domínio restrito, como é o caso do LBot.

A escolha de uma Linguagem de Domínio Específico (DSL), como a proposta LBML, em vez de uma linguagem de propósito geral, tem se mostrado crucial para aumentar a precisão e a interpretabilidade. A DSL restringe o vocabulário e a sintaxe a comandos específicos de movimento, eliminando a ambiguidade inerente à linguagem humana.

Trabalhos recentes têm enfatizado que o uso de uma DSL como intermediário melhora a confiabilidade dos sistemas. A abordagem mediada por DSL (DSL-Mediated Approach), por exemplo, é defendida para evitar a alucinação e a redundância geradas por modelos de linguagem grandes (LLMs) ao traduzir diretamente para linguagens de propósito geral \citep{Wu2025}. Este estudo, focado em estratégias financeiras, demonstrou que a conversão em duas etapas (LN $\rightarrow$ DSL $\rightarrow$ Linguagem de Programação) alcança uma alta taxa de correspondência exata, reforçando o valor da DSL como uma barreira de segurança.

Esta filosofia é implementada diretamente na LBML. Quando o usuário expressa a intenção em português ("LBot, ande 1 metro para a esquerda, gire 45 graus para a direita, depois ande 50 para trás"), o modelo deve mapear essa intenção para a sintaxe rigorosa da LBML (e.g., D100L;R45R;D50B;). O Semantic Parsing para LBML, portanto, não apenas executa uma tarefa, mas garante que o comando gerado seja sintaticamente correto e fisicamente seguro para o robô, um fator crítico em aplicações de robótica.

\subsection{Arquiteturas de Modelos e Mecanismos de Atenção}

O avanço na precisão do Semantic Parsing está intrinsecamente ligado à evolução das arquiteturas de modelos, especialmente os Transformers \citep{Wang2021,Banerji2025}. O trabalho de Banerji et al. \citep{Banerji2025} propôs um framework unificado baseado em um mecanismo de autoatenção com consciência de relações (relation-aware self-attention). Essa técnica representa o esquema como um grafo, onde as arestas codificam relações explícitas (e.g., chave primária), permitindo que o modelo incorpore a estrutura relacional do domínio na sua representação interna.

Essa necessidade de vincular elementos da linguagem natural à estrutura formal se aplica diretamente ao LBot: embora o domínio não seja um banco de dados, o robô opera em um espaço físico com estados e parâmetros. A correta "ligação" (linking) entre substantivos e verbos em português (e.g., "esquerda" ou "virar") e os parâmetros numéricos corretos na LBML é análoga ao Schema Linking tradicional. Portanto, modelos baseados em mecanismos de atenção são ideais para codificar essa estrutura de domínio, mesmo que simplificada para a topologia do robô.

\subsection{Técnicas de Treinamento e Adaptação de Domínio}

A adaptação de modelos de linguagem para domínios específicos é um tema recorrente. As abordagens mais comuns incluem o uso de Large Language Models (LLMs) pré-treinados, seguidos por Fine-Tuning (Ajuste Fino) ou In-Context Learning (ICL). Alternativamente, alguns trabalhos optam por treinar modelos do zero, permitindo maior controle sobre a arquitetura e o comportamento do modelo, embora isso demande mais recursos computacionais e dados de treinamento.

\subsubsection{Fine-Tuning e ICL}

O Fine-Tuning de LLMs (como o proposto no trabalho de NL2SQL para painéis de IoT \citep{Banerji2025}) envolve o uso de um conjunto de dados específico do domínio (e.g., dados de IoT) para adaptar um modelo pré-treinado (e.g., GPT-2) para uma tarefa especializada. Essa abordagem é eficiente em recursos, pois aproveita a base de conhecimento do pré-treinamento. O In-Context Learning (ICL), por sua vez, permite que o modelo aprenda a tarefa apenas por meio de exemplos fornecidos no prompt de entrada, sem a necessidade de atualizar os pesos do modelo \citep{Wu2025}.

\subsubsection{A Escolha do Treinamento do Zero e o Controle da Saída}

Embora a maioria dos trabalhos recentes utilize LLMs pré-treinados, algumas abordagens optam por treinar modelos do zero. O RAT-SQL \citep{Wang2021}, por exemplo, treinou um modelo baseado em relation-aware self-attention especificamente para a tarefa de NL-to-SQL, alcançando resultados competitivos ao capturar as relações estruturais do esquema de banco de dados desde o início do treinamento.

Este trabalho adota uma estratégia semelhante ao optar por treinar um modelo pequeno do zero em um conjunto de dados exclusivo (Português-para-LBML). Esta escolha permite:

\begin{enumerate}
\item \textbf{Evitar dependência de LLMs pré-treinados};
\item \textbf{Garantir controle total sobre o vocabulário e as regras da LBML}.
\end{enumerate}

\subsubsection{Garantia de Estrutura e Sintaxe}

Para garantir que a saída gerada seja sempre um comando LBML sintaticamente válido, o controle da estrutura da DSL é essencial. O conceito de Grammar Prompting \citep{Wang2023} utiliza métricas como a Taxa de Correspondência Exata (Exact Match Rate) em datasets desafiadores como o SPIDER. De forma semelhante, o trabalho de estratégias financeiras \citep{Wu2025} prioriza esta métrica, por ser a única que garante que o comando gerado é 100\% correto. A aplicação dessa métrica é essencial para o LBot, uma vez que, em robótica, qualquer erro na sintaxe da LBML pode resultar em falha de execução ou em um movimento perigoso.

\section{Desenvolvimento da Proposta de Solução}

Esta seção apresenta o planejamento do desenvolvimento do sistema proposto para tradução de comandos em linguagem natural para a LBot Movement Language (LBML). A proposta está organizada em etapas sucessivas, iniciando pela geração do conjunto de dados sintético, seguida da implementação de um modelo base de tradução utilizando arquitetura Transformer, e, posteriormente, do aprimoramento deste modelo com a incorporação do mecanismo de atenção sensível a relações (\textit{Relation-Aware Self-Attention} — RASA). Por fim, será conduzida uma comparação entre as duas versões, de modo a avaliar o impacto da atenção relacional no desempenho do tradutor.

\subsection{Geração do Conjunto de Dados Sintético}

A primeira etapa do projeto consiste na criação de um conjunto de dados sintético para o treinamento inicial do modelo. Pretende-se gerar automaticamente pares de exemplos no formato \textit{Português → LBML}, utilizando um script capaz de combinar diferentes padrões de instruções. Cada amostra representará um comando em linguagem natural e sua respectiva tradução na sintaxe padronizada da LBML.

O dataset abrangerá tanto instruções simples quanto compostas, permitindo ao modelo aprender a estrutura sequencial dos comandos. Serão incluídas variações linguísticas e sinônimos (por exemplo, “ande”, “mova-se”, “siga”) e ruído leve, como variações de acentuação ou ordem dos termos, de modo a aumentar a robustez do modelo diante de entradas reais. Os valores de distância e rotação serão representados em centímetros e graus, respeitando o formato definido pela LBML.

\subsection{Modelo Base sem Relações Explícitas}

A segunda etapa prevê o desenvolvimento de um modelo de tradução baseado na arquitetura Transformer, sem a inclusão inicial de mecanismos de atenção relacional. O modelo foi implementado utilizando uma arquitetura do tipo \textit{Transformer decoder-only}, inspirada em modelos da família GPT. A configuração definida para este trabalho consiste em seis camadas (\textit{layers}), seis cabeças de atenção e vetores de representação de 384 dimensões. O tamanho máximo do contexto foi fixado em 128 tokens, valor adequado para representar comandos contendo até três ações encadeadas.

O vocabulário definido para o modelo possui 70 caracteres, contemplando letras, números, pontuação e os símbolos específicos da LBML. O valor de \textit{dropout} adotado é 0{,}2, buscando garantir estabilidade durante o treinamento sobre um conjunto de aproximadamente 40 mil exemplos sintéticos. Essa configuração foi planejada para oferecer um bom equilíbrio entre capacidade de aprendizado, custo computacional e adequação à complexidade da tarefa de tradução.

O processo de treinamento será acompanhado por avaliações periódicas, utilizando métricas como \textit{Exact Match Rate} e \textit{BLEU Score}, que permitirão monitorar a evolução do modelo e sua capacidade de gerar comandos sintaticamente válidos.

\subsection{Aprimoramento com Atenção Relacional (RASA)}

Após a obtenção de um modelo funcional e estável, será implementada uma segunda versão incorporando o mecanismo de atenção sensível a relações (\textit{Relation-Aware Self-Attention} — RASA). Diferentemente da formulação tradicional, na qual o RASA modela relações entre tokens dentro da própria sequência de entrada, nesta proposta o mecanismo será adaptado para capturar relações entre os tokens da entrada (em português) e os símbolos correspondentes da LBML na saída.

O objetivo é introduzir um viés de atenção que favoreça o alinhamento semântico direto entre os elementos linguísticos e suas representações simbólicas. Por exemplo, o modelo deverá aprender explicitamente relações do tipo “ande” $\rightarrow$ “D”, “gire” $\rightarrow$ “R”, “frente” $\rightarrow$ “F” e “esquerda” $\rightarrow$ “L”. Esse alinhamento relacional entre domínios distintos busca reduzir ambiguidades de tradução e reforçar a consistência sintática dos comandos gerados.

Essas relações serão representadas em uma matriz relacional utilizada durante o processo de \textit{cross-attention}, permitindo que o modelo pondere a influência de cada token de entrada sobre a geração dos símbolos correspondentes da LBML. Espera-se que essa abordagem melhore a correspondência exata em instruções compostas e reduza erros de associação entre ações, direções e parâmetros numéricos.

\subsection{Comparação entre as Versões}

A etapa final do desenvolvimento consistirá na comparação entre o modelo base e o modelo aprimorado com RASA. Ambos serão treinados com o mesmo conjunto de dados e parâmetros idênticos, de modo que as diferenças observadas no desempenho possam ser atribuídas exclusivamente à introdução da atenção relacional.

A métrica de avaliação será a \textit{Exact Match Rate}. Serão analisadas separadamente instruções simples e compostas, a fim de verificar se o ganho de desempenho é mais significativo em cenários com maior complexidade sintática.

% Bibliografia
\bibliographystyle{sbc}
\bibliography{sbc-template}

\end{document}